\chapter{Scientific and theoretical background}
\section{Amyloid-β40 and intrinsically disordered ensembles}
An ensemble describes the conformations a molecule can adopt and the statistical weight of each one. For an intrinsically disordered protein (IDP), a representative structure shows only one member of the ensemble, whereas an ensemble-averaged measurement reflects many conformations. Global descriptors summarize this information but lose some structural detail. Two configurations can have the same radius of gyration and solvent-accessible area but different residue contacts. Clustering these descriptors groups conformations by the measured similarities; it does not identify a unique atomic fold.

Experimental data can help constrain an ensemble, although NMR and SAXS measurements generally average over conformations. Borthakur and colleagues describe a maximum-entropy procedure that combines simulations with experimental data to constrain IDP ensembles \autocite{borthakur2025}. Their study offers a possible independent comparison for this project, provided that the structural data, measured quantities and statistical weights are compatible.

\section{Molecular dynamics and conformational ensembles}
Molecular dynamics generates configurations by numerically integrating the equations of motion under a specified potential-energy function. Force-field parameters, the solvent model, thermodynamic conditions and sampling all affect the resulting ensemble. The primary study used GROMACS, CHARMM36m and PLUMED \autocite{cadenelli2026,abraham2015,huang2017,tribello2014}. This project starts from the descriptor table produced after those simulations.

In an equilibrium canonical ensemble, an observable $A(q)$ has expectation proportional to the integral of $A(q)\exp[-\beta_T U(q)]$, where $q$ denotes coordinates and $U$ is potential energy. For energies per molecule, $\beta_T=1/(k_B T)$; for the molar energies used in the implementation, the consistent factor is $1/(RT)$. Mixing these energy conventions introduces a scale error even when the reweighting code is numerically stable.

\section{Sampling limitations}
Saving configurations more frequently adds rows to a trajectory without necessarily adding independent information. Nearby frames may remain in the same slowly changing region of descriptor space. Under suitable assumptions, statistical inefficiency measures how this dependence increases the variance of an estimate. It cannot account for a region of the ensemble that was never sampled. Even a stable overall mean can conceal changes in the relative populations of different structural groups.

Assessing equilibration and sampling quality requires more than a raw frame count \autocite{chodera2016,grossfield2019}. Here, the trajectory was also generated under a changing bias. Correlation statistics describe dependence in that recorded process; they do not give unbiased relaxation times, transition rates, or kinetic states. The reproduction retains the published analysis interval. The extension examines the statistics within that interval without changing its boundaries.

\section{Well-tempered metadynamics}
WT-MetaD adds a history-dependent potential along selected CVs. Its tempering parameter is conventionally expressed as $\gamma=(T+\Delta T)/T>1$. In the long-time regime, the bias approaches a scaled negative free energy, up to an additive constant \autocite{barducci2008}. This relationship motivates the final-bias approximation used in the primary paper.

The bias encourages exploration but also changes how often conformations are sampled and the forces acting on the peptide. Population estimates therefore need a statistical correction, and the recorded time cannot be read directly as an unbiased physical clock. Thermodynamic reweighting corrects the distribution under stated assumptions; recovering kinetic information is a separate problem.

\section{Gaussian hills and bias construction}
For CV vector $s$, hills deposited at times $t_k$ generate a bias of the form
\begin{equation}
V(s,t)=\sum_{t_k<t}h_k\exp\!\left[-\frac12\sum_a\frac{(s_a-s_a(t_k))^2}{\sigma_a^2}\right].
\label{eq:hills}
\end{equation}
In WT-MetaD, hill heights decrease with the accumulated bias as $h_k=h_0\exp[-V(s(t_k),t_k)/(k_B\Delta T)]$, using molecular energy units. For molar units, $k_B$ is replaced by $R$. The convention must agree with the simulation software and energy units. Equation~\ref{eq:hills} explains how the bias is constructed. It has not been reconstructed here because no project HILLS file is available.

PLUMED distinguishes the instantaneous bias output, the time-dependent offset, and their difference. Its \code{metad.rbias} component equals \code{metad.bias} minus \code{metad.rct}; the latter represents $c(t)$ \autocite{plumedmetad}. A column named \code{metad.rbias} is therefore potentially useful for a separate time-dependent reweighting analysis. Its name alone does not establish that the offset was computed with the required settings or that its units are known.

\section{Free-energy surfaces}
For a CV distribution $p(s)$, free energy is $F(s)=-\beta_T^{-1}\log p(s)+C$, where $C$ sets an arbitrary reference. This surface describes the chosen CVs. Integrating over the remaining coordinates can hide structural differences, so a two-dimensional minimum does not uniquely specify a three-dimensional conformation.

PLUMED's \code{sum\_hills} supports reconstruction from accumulated hill records and distinguishes free-energy output from negative-bias output \autocite{plumedsum}. That distinction matters because the WT scaling factor must not be applied twice. The recovered complete file labels its energy field \code{file.free}, but the energy unit, generating commands, and replica-combination procedure are still needed for a verified final-bias calculation.

\section{Statistical reweighting}
For a quasi-static bias, the correction from a biased distribution to the corresponding unbiased distribution has weights proportional to $\exp(\beta_T V)$. Combining this with the WT relationship yields
\begin{equation}
\ell_i\equiv\log w_i=-\frac{1-1/\gamma}{RT}F(\alpha_i,\beta_i).
\label{eq:finalweights}
\end{equation}
The free-energy zero does not affect the normalized result: adding a constant to $F$ adds the same offset to all relevant log weights, which cancels on normalization. Equation~\ref{eq:finalweights} implements the primary paper's final-bias approximation. It assumes that the analyzed interval can be treated as quasi-stationary, with a sufficiently stable bias for this correction to apply.

Alternative MetaD reweighting approaches account explicitly for changing bias normalization \autocite{bonomi2009,tiwary2015}. For verified PLUMED offset-corrected output, the optional comparison uses $\ell_i=\mathrm{rbias}_i/(RT)$ \autocite{plumedreweight}. These two estimators are not interchangeable labels for the same computation. The current repository keeps their configurations and resulting outputs separate.

\section{Structural descriptors}
The baseline uses $\na$, $\nb$, $R_g$, and total SASA. In the primary study, the first two are scaled continuous secondary-structure CVs, intended to approximate residue counts. They should not be interpreted as exact integer DSSP assignments. DSSP itself classifies local structural patterns using geometrical and hydrogen-bond information \autocite{kabsch1983}; reproducing such assignments requires coordinates, which are absent here.

Radius of gyration summarizes global compaction; SASA measures solvent exposure. They may be correlated while still distinguishing different conformations within a cluster. Apolar and polar SASA, intrapeptide hydrogen bonds, and peptide--solvent hydrogen bonds provide further descriptors after baseline clustering. Their exact definitions must be recovered from the original analysis settings, including whether the two SASA components add up to total SASA.

\section{Density-based conformational clustering}
Density-peak methods seek centers that have both high local density and a large distance from denser points. The primary study applies this idea repeatedly to find acceptable clusters across random subsets and density scales \autocite{cadenelli2026}. A silhouette coefficient compares how close a point is to its own cluster with how close it is to other clusters \autocite{rousseeuw1987}. It measures geometric separation and does not establish thermodynamic metastability.

Consensus then groups local clusters by similarities in their descriptor distributions. DPA offers a distinct, statistically motivated density-topography approach \autocite{derrico2021}. A DPA benchmark would test whether an alternative treatment of density produces similar families. Likewise, weighted ShapeGMM addresses structural ensembles in a coordinate-based representation \autocite{sasmal2024}; its applicability here depends on obtaining valid structural frames.

\begin{figure}[htbp]
\centering
\begin{subfigure}{.485\textwidth}
\includegraphics[width=\linewidth]{artwork/ensemble-editorial-v2.png}
\caption{Conformational diversity}
\end{subfigure}\hfill
\begin{subfigure}{.485\textwidth}
\includegraphics[width=\linewidth]{artwork/bias-editorial-v2.png}
\caption{An evolving sampling landscape}
\end{subfigure}
\caption{Conceptual editorial illustrations. Open paper ribbons evoke different conformations of a flexible peptide; layered deposits evoke the accumulation of metadynamics bias. Neither image depicts molecular coordinates, secondary-structure assignments, a measured FES or a physical trajectory. The qualitative ideas are formalized in the surrounding text and \cref{eq:hills}. Generated-art provenance is documented in Appendix D.}
\label{fig:editorialconcepts}
\end{figure}
